\chapter{成績の測り方}
\label{ch:perf}

戦略の良し悪しを比べるには、共通の物差しが要る。この章では、運用の世界で使われる成績の指標と、その落とし穴をまとめる。

\section{リターン}

\subsection{単純リターンと対数リターン}

時点 $t-1$ から $t$ までの資産の\textbf{単純リターン}は、価格 $P_t$ と配当 $D_t$ から
\begin{equation}
R_t = \frac{P_t + D_t}{P_{t-1}} - 1
\end{equation}
である。\textbf{対数リターン}は $r_t = \ln(1+R_t)$。対数リターンは期間をまたいで足し算できる（$r_{1\to T} = \sum_t r_t$）。一方、ポートフォリオのリターンは単純リターンの加重平均（$R_{p,t} = \sum_i w_i R_{i,t}$）で、対数リターンでは成り立たない。銘柄をまたぐ計算には単純リターン、時間をまたぐ計算には対数リターンが便利である。

\subsection{超過リターンと年率化}

\begin{itemize}
\item \textbf{超過リターン}：安全な資産（短期国債など）の金利 $R_f$ を引いた $R_t - R_f$。ヘッジファンドは現金を担保に先物や空売りを使うので、比べる相手は「現金」であることが多い。
\item \textbf{比べる相手との差}（アクティブ・リターン）：指数に対して運用するなら、$R_{p,t} - R_{b,t}$。
\item \textbf{年率化}：日次の平均 $\bar R$ と標準偏差 $\sigma$ から、年率の平均は $252\,\bar R$、年率の標準偏差は $\sqrt{252}\,\sigma$（月次なら 12 と $\sqrt{12}$）。標準偏差に平方根がつくのは、独立な変動の分散が期間に比例して足されるからである。
\end{itemize}

\subsection{算術平均と幾何平均}

年ごとのリターンの算術平均を $\mu$、標準偏差を $\sigma$ とすると、複利で増える速さ（幾何平均）$g$ はおよそ
\begin{equation}
g \approx \mu - \frac{\sigma^2}{2}
\end{equation}
である。同じ平均でも、変動が大きいほど資産は増えにくい（ボラティリティの引きずり）。例えば $+50\%$ と $-50\%$ を交互に繰り返すと、算術平均は 0 だが、資産は 2 年で $1.5 \times 0.5 = 0.75$ 倍になる。レバレッジを上げると $\mu$ は比例して増えるが、$\sigma^2$ は 2 乗で増えるので、ある倍率を超えると $g$ はかえって下がる（\chapref{ch:pc} のケリー基準）。

\section{シャープレシオ}

\subsection{定義}

最も広く使われる物差しが、Sharpe（1966）\cite{sharpe1966} の\textbf{シャープレシオ}である。
\begin{equation}
\SR = \frac{\E[R - R_f]}{\sigma(R - R_f)}
\end{equation}
リスク 1 単位あたりの超過リターンを表す。レバレッジを $k$ 倍にすると、分子も分母も $k$ 倍になるので、シャープレシオは変わらない。つまり、シャープレシオの高い戦略は、借入れで好きなリスクの水準に調整できる。ヘッジファンドがシャープレシオを重視するのは、このためである。

日次のシャープレシオ $\SR_d$ を年率にするには、リターンが独立なら $\SR = \sqrt{252}\,\SR_d$ とする。

\begin{table}[H]
\centering\small
\caption{シャープレシオの目安（年率、費用の後）}
\label{tab:sr}
\begin{tabular}{@{}lP{0.66\linewidth}@{}}
\toprule
年率の SR & 目安 \\
\midrule
0.3〜0.5 & 米国株式の長期の平均（市場全体をただ持つ） \\
0.5〜1.0 & よく分散したファクター戦略、トレンドフォローの長期の成績 \\
1.0〜2.0 & 優れたヘッジファンド。長く続くことは少ない \\
2 を超える & 容量の小さい短期の戦略、マーケットメイク。あるいはバックテストの過学習 \\
\bottomrule
\end{tabular}
\end{table}

\subsection{推定の誤差}

シャープレシオは推定値であり、誤差が大きい。リターンが独立で正規分布に従うとき、$T$ 期の標本から推定したシャープレシオ $\widehat{\SR}$（同じ期間の単位）の標準誤差は、Lo（2002）\cite{lo2002} によればおよそ
\begin{equation}
\mathrm{SE}(\widehat{\SR}) \approx \sqrt{\frac{1 + \SR^2/2}{T}}
\end{equation}
である。年率で考えると、$\SR$ が小さいとき、$t$ 値はおよそ $\SR \times \sqrt{\text{年数}}$ になる。

\begin{exbox}[腕を見分けるのに何年かかるか]
年率のシャープレシオが $\SR$ の戦略が、偶然でないと言えるには（$t \geq 2$）、およそ $(2/\SR)^2$ 年の記録が要る。
\begin{center}
\begin{tabular}{@{}cccccc@{}}
\toprule
年率の SR & 0.25 & 0.5 & 1.0 & 2.0 & 3.0 \\
\midrule
必要な年数 & 64 & 16 & 4 & 1 & 0.44 \\
\bottomrule
\end{tabular}
\end{center}
年率 0.5 の戦略（よいファクター戦略）が本物かどうかを、成績の記録だけで確かめるには 16 年かかる。しかも、その間に戦略の中身や市場が変わる。短期の戦略が好まれるもう一つの理由は、同じ年数で多くの賭けができ、早く検証できることである。
\end{exbox}

\subsection{シャープレシオの限界}

シャープレシオは、平均と標準偏差しか見ない。次のような戦略は、シャープレシオが実際より良く見える。

\begin{itemize}
\item \textbf{オプションの売り}：普段は小さな利益を積み重ね、まれに大きく損をする（負の歪度）。損の起きない期間を切り取ると、シャープレシオは非常に高い。
\item \textbf{評価の平滑化}：市場の価格がない資産を、運用者の評価や古い価格で値付けすると、リターンの変動が小さく見え、自己相関が出る（\ref{sec:smoothing} 節）。
\item \textbf{レバレッジの解消の連鎖}：同じ取引をしている人が一斉に売ると、普段の標準偏差からは想像できない損失が出る（第 II 部、第 \ref{ch:failures} 章）。
\end{itemize}

そこで、次の指標も一緒に見る。

\section{そのほかの指標}

\begin{description}[leftmargin=1.5em,style=nextline]
\item[ソルティノ・レシオ]
  分母を、下方の変動（目標を下回ったリターンだけの標準偏差）にしたもの。上方の変動を罰しない。
\item[最大ドローダウン]
  資産の過去の最高値からの最大の下落率。
  \begin{equation}
  \mathrm{MDD} = \max_t \frac{\max_{s \le t} V_s - V_t}{\max_{s \le t} V_s}
  \end{equation}
  投資家が実際に耐えられるかどうかを最もよく表す。期間が長いほど大きくなるので、同じ期間で比べる。
\item[カルマー・レシオ]
  年率のリターンを最大ドローダウンで割ったもの。
\item[歪度と尖度]
  リターンの分布の非対称と、裾の厚さ。負の歪度と厚い裾は、まれな大損を示す。
\item[VaR と期待ショートフォール]
  損失の分布の分位点と、その先の平均（\chapref{ch:risk}）。
\item[回転率]
  1 年に売買する額を資産で割ったもの。高いほど取引コストがかかる（\chapref{ch:cost}）。
\end{description}

\section{ベータ・アルファ・情報レシオ}

\subsection{回帰による分解}

戦略の超過リターンを、市場の超過リターンに回帰する。
\begin{equation}
R_{p,t} - R_{f,t} = \alpha + \beta\,(R_{m,t} - R_{f,t}) + \varepsilon_t
\end{equation}
\begin{itemize}
\item $\beta$：市場が 1\% 動いたときに戦略が動く大きさ。指数のファンドで安く手に入る部分。
\item $\alpha$：市場で説明できない平均の超過リターン。「腕」の候補。
\item $\varepsilon_t$：市場と無関係な変動。
\end{itemize}
市場だけでなく、規模・割安・モメンタムなどのファクター（\chapref{ch:factor}）も右辺に入れると、「ファクターでは説明できないアルファ」が測れる。ヘッジファンドの研究では、オプションに似たトレンドフォローのファクターなども使う（第 II 部、第 \ref{ch:hfperf} 章）。

\subsection{トラッキングエラーと情報レシオ}

指数に対して運用するとき、比べる相手との差の標準偏差を\textbf{トラッキングエラー}（TE）と呼ぶ。
\begin{equation}
\TE = \sigma(R_p - R_b), \qquad \IR = \frac{\E[R_p - R_b]}{\TE}
\end{equation}
$\IR$ は\textbf{情報レシオ}で、比べる相手との差についてのシャープレシオである。TE の上限は、「指数からどれだけ離れてよいか」の予算になる。LLM-BLM のポートフォリオで年 3\% の TE の上限を置いたのは、この考え方である。

\section{アクティブ運用の基本法則}
\label{sec:flam}

Grinold（1989）\cite{grinold1989} は、情報レシオを、予測の質と賭けの数に分解した。
\begin{equation}
\IR \approx \IC \times \sqrt{\mathrm{BR}}
\label{eq:flam}
\end{equation}
\begin{itemize}
\item $\IC$（information coefficient、情報係数）：予測と、実際の（リスクで割った）リターンの相関。0.05 なら「よい予測」、0.1 を超えることは珍しい。
\item $\mathrm{BR}$（breadth、幅）：1 年に行う独立な賭けの数。500 銘柄を毎月入れ替えるなら、最大で $500 \times 12 = 6{,}000$。
\end{itemize}

Clarke・de Silva・Thorley（2002）\cite{cdt2002} は、制約（ロングのみ、TE の上限など）で予測を比率に移しきれない分を、\textbf{移転係数} $\mathrm{TC}$（0〜1）として加えた。
\begin{equation}
\IR \approx \mathrm{TC} \times \IC \times \sqrt{\mathrm{BR}}
\end{equation}

\begin{exbox}[基本法則で考える]
\begin{itemize}
\item IC 0.05、500 銘柄、毎月独立に予測：$\IR = 0.05 \times \sqrt{6000} \approx 3.9$。
\item 同じで、ロングのみの制約で $\mathrm{TC} = 0.3$ だったとする（仮の値）：$\IR \approx 1.2$。
\item さらに、銘柄の予測が互いに相関していて、実質的に独立な賭けが 1 年に 600 回しかない：$\IR \approx 0.3 \times 0.05 \times \sqrt{600} \approx 0.37$。
\item 取引コストを引くと、さらに下がる。
\end{itemize}
この式は、弱い予測でも多くの銘柄・多くの時点に広げれば価値があることを示す一方、独立性・制約・コストで、その価値の大半が失われうることも示す。クオンツ運用が数千銘柄を扱い、制約を減らし（空売り、ロング・ショート）、コストを細かくモデル化するのは、この式の各項を大きくするためである。
\end{exbox}

\section{評価の平滑化と自己相関}
\label{sec:smoothing}

流動性の低い資産（社債、私募の資産、小型株）は、毎月の市場の価格が信頼できない。運用者が古い価格や推定で値付けすると、報告されるリターンは、本当のリターンを数か月にわたってならしたものになる。Getmansky・Lo・Makarov（2004）\cite{glm2004} は、報告されるリターン $R^o_t$ を
\begin{equation}
R^o_t = \theta_0 R_t + \theta_1 R_{t-1} + \cdots + \theta_k R_{t-k}, \qquad \textstyle\sum_j \theta_j = 1
\end{equation}
とモデル化した。ならすと、平均は変わらないが、標準偏差は小さくなり、シャープレシオは高く見え、リターンに正の自己相関が出る。

自己相関があるときの年率化は、$\sqrt{12}$ 倍では正しくない。月次の自己相関を $\rho_k$ とすると、$q$ か月の和の分散は
\begin{equation}
\Var\Big(\sum_{j=1}^{q} R_j\Big) = \sigma^2\Big(q + 2\sum_{k=1}^{q-1}(q-k)\rho_k\Big)
\end{equation}
で、正の自己相関があると、年率の変動は $\sqrt{12}\,\sigma$ より大きい。Lo（2002）は、この補正をしたシャープレシオを示した。

\section{手数料の前と後}

ヘッジファンドの成績は、手数料の前（グロス）と後（ネット）で大きく違う（\ref{sec:fees} 節の例）。論文や資料のシャープレシオが、手数料・取引コスト・借入れの費用・貸株料のどれを引いた後の値かを、必ず確かめる必要がある。研究の戦略の成績は、多くが手数料も取引コストも引く前の値である。

\begin{projbox}
LLM-BLM の月次の運用で、年率の TE を 3\% とし、情報レシオが 0.3 だったとすると、比べる相手との差の期待値は年 0.9\%。$t \geq 2$ に必要な年数は $(2/0.3)^2 \approx 44$ 年である。数年の検証期間で「有意に勝った」と言える戦略は、情報レシオが 1 前後以上でなければならない。この研究で時価総額加重との差が有意にならなかったことは、それだけでは「予測が無価値」を意味しない。ただし、2 日の反応で測った IC と、月次の運用で実現した IC の大きな差（情報が当日に株価に入りきる）は、検定力とは別の問題である。
\end{projbox}
